\Subhead{1.3.2}{Гашение колебаний}
\textbf{Гашение} --- подавление самостоятельной составляющей колебаний сильным внешним сигналом. Активный элемент насыщается, и усиления для свободных колебаний уже не хватает. Движение не прекращается: остаётся вынужденный ответ на частоте источника. При захвате подчиняется фаза уже существующей генерации; при гашении подавляется её самостоятельный механизм.

Это видно в модели Ван-дер-Поля. Вне близкого резонанса положим $x\simeq b(t)\cos\omega_gt+B\cos\Omega t$, где $B$ --- вынужденная, $b$ --- свободная амплитуда. Усредняя отдельно по двум фазам при отсутствии низких комбинационных резонансов, получаем
\begin{equation}\label{eq:quenching}
 \dot b=\mu b\left(1-\frac{b^2}{4a^2}-\frac{B^2}{2a^2}\right).
\end{equation}
Член $-B^2/(2a^2)$ --- перекрёстное насыщение. При $B^2>2a^2$ малые свободные колебания затухают. Это модельный порог, а не универсальная характеристика транзистора. При близких частотах нельзя независимо усреднить обе фазы: нужны полные уравнения \eqref{eq:forcedSL}.
% Самостоятельный вывод: резонансная проекция x^2 xdot для двух несоизмеримых частот.
